
Zur weiteren Analyse der Schaltung wird zusätzlich die am Lastwiderstand anliegende
Spannung normiert dargestellt. Dadurch lässt sich der Spannungsverlauf unabhängig
von der absoluten Höhe der Quellenspannung untersuchen. Die relative Spannung wird
definiert als:
\begin{equation}
u = \frac{U}{U_0}
\end{equation}

\noindent
Die am Lastwiderstand anliegende Spannung $U$ ist durch den Spannungsteiler aus
Gleichung (2) gegeben. Durch Einsetzen dieser Beziehung in die Definition der rela-
tiven Spannung ergibt sich:
\begin{equation}
u = \frac{R_L}{R_I + R_L}
\end{equation}

\noindent
Zur Darstellung in normierter Form wird erneut das Widerstandsverhältnis
\begin{equation}
r = \frac{R_L}{R_I}
\end{equation}
verwendet. Damit ergibt sich für die relative Spannung die folgende Abhängigkeit vom
Widerstandsverhältnis:
\begin{equation}
u(r) = \frac{r}{1 + r}
\end{equation}

\begin{figure}[H]
    \centering
    \includegraphics[width=1\textwidth]{2_1_4_u_r_spannungsquelle.png}
    \caption{Spannung im Verhältnis der Widerstände (Spannungsquelle)}
    \label{fig:Spannung im Verhaeltnis der Widerstaende (Spannungsquelle)}
\end{figure}

\noindent
Die relative Spannung nimmt mit wachsendem Lastwiderstand monoton zu. Für sehr
kleine Lastwiderstände ($r \ll 1$) geht die relative Spannung gegen null, während sie
für große Lastwiderstände ($r \gg 1$) gegen eins strebt. Die Darstellung erfolgt mit
logarithmischer Skalierung der $r$-Achse, wodurch der Spannungsverlauf über mehrere
Größenordnungen anschaulich dargestellt werden kann.
